% ===== PRÉAMBULE COMMUN — À COPIER TEL QUEL EN TÊTE DE CHAQUE FICHIER .tex =====
% Ne modifier QUE :
%   - les deux lignes \fancyhead[L] et \fancyhead[R]
%   - le bloc titre \begin{center}...\end{center} juste après \begin{document}
\documentclass[11pt,a4paper]{article}
\usepackage[utf8]{inputenc}
\usepackage[T1]{fontenc}
\usepackage{lmodern}
\usepackage[french]{babel}
\usepackage{amsmath,amssymb,mathtools}
\usepackage{icomma}
\usepackage{xcolor}
\usepackage{tikz}
\usepackage{pgfplots}
\usepackage[most]{tcolorbox}
\usepackage{enumitem}
\usepackage{array,booktabs,colortbl}
\usepackage[a4paper,margin=2.2cm,headheight=26pt,headsep=10pt]{geometry}
\usepackage{fancyhdr}
\usepackage[hidelinks]{hyperref}
\usetikzlibrary{arrows.meta,calc,angles,quotes,positioning,decorations.markings,intersections,backgrounds}
\pgfplotsset{compat=1.18}

\definecolor{primary}{RGB}{226,74,88}
\definecolor{primarydark}{RGB}{150,28,42}
\definecolor{methcol}{RGB}{40,90,150}
\definecolor{excol}{RGB}{0,135,90}
\definecolor{defcol}{RGB}{0,114,178}
\definecolor{attncol}{RGB}{200,40,55}
\definecolor{thmcol}{RGB}{0,140,120}
\definecolor{courbe}{RGB}{32,110,200}
\definecolor{courbeder}{RGB}{210,70,40}
\definecolor{tang}{RGB}{0,150,90}
\definecolor{grille}{RGB}{200,205,215}
\definecolor{plus}{RGB}{200,60,60}
\definecolor{moins}{RGB}{30,90,200}

\tcbset{boxbase/.style={enhanced,breakable,boxrule=0.9pt,arc=2.5pt,
  left=3mm,right=3mm,top=2.3mm,bottom=2.3mm,fonttitle=\bfseries,coltitle=white}}
\newtcolorbox{cle}[1][]{boxbase,colback=defcol!6,colframe=defcol,
  title={\textsc{À retenir}\ifx\relax#1\relax\relax\else\textmd{ : \itshape #1}\fi}}
\newtcolorbox{methode}[1][]{boxbase,colback=methcol!7,colframe=methcol,sharp corners,
  title={\textsc{Méthode}\ifx\relax#1\relax\relax\else\textmd{ --- \itshape #1}\fi}}
\newtcolorbox{exemplebox}[1][]{boxbase,colback=excol!5,colframe=excol,
  title={\textsc{Exemple}\ifx\relax#1\relax\relax\else\textmd{ : \itshape #1}\fi}}
\newtcolorbox{piege}[1][]{boxbase,colback=attncol!6,colframe=attncol,
  title={\textsc{Attention}\ifx\relax#1\relax\relax\else\textmd{ : \itshape #1}\fi}}
\newtcolorbox{exo}[1][]{boxbase,colback=primary!4,colframe=primary,
  title={\textsc{Exercice}\ifx\relax#1\relax\relax\else\textmd{~#1}\fi}}
\newtcolorbox{corr}[1][]{boxbase,colback=thmcol!5,colframe=thmcol,
  title={\textsc{Corrigé}\ifx\relax#1\relax\relax\else\textmd{~#1}\fi}}

\newcommand{\R}{\mathbb{R}}
\newcommand{\N}{\mathbb{N}}
\newcommand{\ds}{\displaystyle}
\newcommand{\tg}{\operatorname{tg}}
\newcommand{\cotg}{\operatorname{cotg}}
\newcommand{\arctg}{\operatorname{arctg}}
\newcommand{\arccotg}{\operatorname{arccotg}}
\newcommand{\limh}{\lim_{h\to 0}}

\pagestyle{fancy}\fancyhf{}
\fancyhead[L]{\small\textcolor{primarydark}{\textbf{Thème 1} — Nombre dérivé \& taux d'accroissement}}
\fancyhead[R]{\small\textcolor{primarydark}{\textsc{Corrigé — Moyen}}}
\fancyfoot[C]{\small\thepage}
\renewcommand{\headrulewidth}{0.8pt}

\begin{document}

\begin{center}
{\LARGE\bfseries\color{primarydark} Corrigé — Niveau Moyen}\\[2pt]
{\color{primary}\rule{0.5\textwidth}{1.2pt}}\\[4pt]
{\small\itshape Nombre dérivé $f'(x_0)$ obtenu par passage à la limite}
\end{center}
\vspace{2mm}

\begin{corr}[1 --- $f(x)=x^2$, montrer $f'(x_0)=2x_0$]
\textbf{Taux d'accroissement :}
$\ds \tau(h)=\frac{(x_0+h)^2-x_0^2}{h}$.\\[2pt]
\textbf{Numérateur :} $(x_0+h)^2-x_0^2=x_0^2+2x_0h+h^2-x_0^2=2x_0h+h^2=h(2x_0+h)$.\\[2pt]
\textbf{Simplification par $h$ :} $\ds \tau(h)=\frac{h(2x_0+h)}{h}=2x_0+h$.\\[2pt]
\textbf{Limite :} $\ds f'(x_0)=\limh(2x_0+h)=\boxed{2x_0}$. \hfill$\square$
\end{corr}

\begin{corr}[2 --- $f(x)=x^2+3x$]
\[
\tau(h)=\frac{\big[(x_0+h)^2+3(x_0+h)\big]-\big[x_0^2+3x_0\big]}{h}.
\]
Numérateur : $x_0^2+2x_0h+h^2+3x_0+3h-x_0^2-3x_0=2x_0h+h^2+3h=h(2x_0+h+3)$. Donc
\[
\tau(h)=2x_0+h+3\ \xrightarrow[h\to 0]{}\ \boxed{2x_0+3}.
\]
\end{corr}

\begin{corr}[3 --- $f(x)=x^3$]
\[
\tau(h)=\frac{(x_0+h)^3-x_0^3}{h}.
\]
On développe : $(x_0+h)^3=x_0^3+3x_0^2h+3x_0h^2+h^3$, donc le numérateur vaut
\[
3x_0^2h+3x_0h^2+h^3=h\big(3x_0^2+3x_0h+h^2\big).
\]
D'où $\ds \tau(h)=3x_0^2+3x_0h+h^2$, et en faisant $h\to 0$ :
\[
f'(x_0)=\boxed{3x_0^2}.
\]
\end{corr}

\begin{corr}[4 --- $f(x)=\tfrac{1}{x}$, $x_0\neq 0$]
\[
\tau(h)=\frac{\dfrac{1}{x_0+h}-\dfrac{1}{x_0}}{h}.
\]
On réduit le numérateur au même dénominateur $x_0(x_0+h)$ :
\[
\frac{1}{x_0+h}-\frac{1}{x_0}=\frac{x_0-(x_0+h)}{x_0(x_0+h)}=\frac{-h}{x_0(x_0+h)}.
\]
Donc
\[
\tau(h)=\frac{1}{h}\cdot\frac{-h}{x_0(x_0+h)}=\frac{-1}{x_0(x_0+h)}
\ \xrightarrow[h\to 0]{}\ \frac{-1}{x_0\cdot x_0}=\boxed{-\frac{1}{x_0^{2}}}.
\]
\end{corr}

\begin{corr}[5 --- mobile MRUA, $y=f(x)=2x^2$]
\textbf{Vitesse instantanée par la définition :}
\[
\tau(h)=\frac{2(x_0+h)^2-2x_0^2}{h}
=\frac{2\big(x_0^2+2x_0h+h^2\big)-2x_0^2}{h}
=\frac{4x_0h+2h^2}{h}=4x_0+2h.
\]
En faisant $h\to 0$ : $\ds f'(x_0)=\boxed{4x_0}$.\\[3pt]
\textbf{À l'instant $t=3$~s :} $f'(3)=4\times 3=\boxed{12\ \text{m/s}}$.
\end{corr}

\end{document}
